% !TEX program = xelatex
\documentclass[11pt]{article}
\usepackage[a4paper,margin=24mm]{geometry}
\usepackage{amsmath,amssymb,amsthm,mathtools}
\usepackage{xcolor}
\usepackage[colorlinks=true,linkcolor=blue!50!black,citecolor=blue!50!black,urlcolor=blue!50!black]{hyperref}
\hypersetup{pdftitle={The locally nilpotent part of nonboundary localization quotients (corrected)},pdfauthor={Discussion note}}
\newcommand{\C}{\mathbb C}
\newcommand{\Z}{\mathbb Z}
\newcommand{\g}{\widehat{\mathfrak{sl}}_2}
\newcommand{\slc}{\mathfrak{sl}_2}
\newcommand{\ad}{\operatorname{ad}}
\newcommand{\im}{\operatorname{Im}}
\newcommand{\D}{\mathcal D}
\newcommand{\Span}{\operatorname{span}_{\C}}
\newtheorem{theorem}{Theorem}
\newtheorem{lemma}[theorem]{Lemma}
\newtheorem{proposition}[theorem]{Proposition}
\newtheorem{corollary}[theorem]{Corollary}
\setlength{\emergencystretch}{2em}
\title{The Locally Nilpotent Part of Nonboundary Localization Quotients\\
  \large Corrected E-action and primitive recursion; the weight-filtration
  claims are withdrawn}
\author{Progress report, corrected after the audit of 2026-09-22}
\date{September 22, 2026}
\begin{document}
\maketitle

\begin{abstract}
This is a corrected version of the September 22 note on
$G_c=\{v\in\mathcal T_{f_c}M_{a,b}:e_{-c}^nv=0$ for some $n\}$ in the deep
region. The independent audit (\texttt{audits/2026-09-22/}) found explicit
counterexamples to the E-action formula, to the primitive-space description
$P_\nu\cong\{X:T^{\nu+1}(X)u=0\}$, to the claim $P_0=\C w_1$ at $\mu=-2$,
and to the statement that the weight tails $G^{\ge\nu}$ are affine
submodules. This version records the corrected E-action through the full
components $A_r$, the corrected triangular recursion for primitive vectors,
the infinite family $Z^nw_1$ in $P_0$, the complete set of projections
including the lowering map $\ell_j:P_\nu\to P_{\nu-2}$, and withdraws the
weight-filtration consequences. The statements proved earlier and not
affected are the boundary isomorphism, the deep-region simplicity criterion,
the unique maximal submodule with simple head, and the algebraic central
reduction compatibility. The internal classification of $G_c$ (simplicity,
socle, length) is open.
\end{abstract}

\section{Definitions, notation, and assumptions}

\paragraph{Affine algebra and induced modules.}
As in the unified note: $\g=(\mathfrak{sl}_2\otimes\C[t,t^{-1}])\oplus\C K$,
\[
 [e_i,f_j]=h_{i+j}+i\delta_{i+j,0}K,\quad [h_i,h_j]=2i\delta_{i+j,0}K,\quad
 [h_i,e_j]=2e_{i+j},\quad [h_i,f_j]=-2f_{i+j},
\]
$K$ central, $[e_i,e_j]=[f_i,f_j]=0$. Fix $a,b\in\Z$, $a+b\ge0$,
$L=a+b+1\ge1$, and
$S_{a,b}=\Span\{K,h_i\ (i\ge0),e_i\ (i>a),f_j\ (j>b)\}$. The character
$\chi:S_{a,b}\to\C$ is determined by $\kappa=\chi(K)$ and
$\lambda_i=\chi(h_i)$, $0\le i\le L$, with $\lambda_i=0$ for $i>L$;
$M_{a,b}(\chi)=U\otimes_{U(S_{a,b})}\C_\chi$, $u=1\otimes1_\chi$. The central
level is denoted by $\kappa$ (not $c$), reserving $c$ for the localization
index.

\paragraph{Deep-region localization.}
Fix $c<b$ with $c\le-a-1$ (equivalently $d_*:=b-c\ge L$) and set $F=f_c$.
The spectral-flow twist $\sigma_{-c}$ identifies $\mathcal T_FM_{a,b}(\chi)$
with $Q=\mathcal T_{f_0}M_{A,B}(\chi')$, where $A=a+c\le-1$, $B=d_*\ge L$,
and $\mu:=\chi'(h_0)=\lambda_0-c\kappa$. Write $E=e_0$, $H=h_0$,
$w_m=f_0^{-m}u+M$. Then $[E,f_0]=H$, $Eu=0$, $Hu=\mu u$.

\paragraph{PBW notation.}
\[
 \mathcal X=\Big\{X=f_{j_1}\cdots f_{j_p}e_{i_1}\cdots e_{i_q}h_{r_1}\cdots
 h_{r_s}:j_t\le B,\ j_t\ne0,\ i_t\le A,\ r_u<0\Big\},
\]
including the empty word. $\{F^{-m}Xu:m\ge1,\ X\in\mathcal X\}$ is a basis
of $Q$. Write $\deg X=2q-2p$; then $H(F^{-m}Xu)=(\mu+2m+\deg X)F^{-m}Xu$.
Let $V_0\subset M$ be the span of $\{Xu:X\in\mathcal X\}$ (the $F$-free
part), so that $M=\bigoplus_{r\ge0}f_0^rV_0$.

\paragraph{The locally nilpotent part.}
$G:=\{v\in Q:E^nv=0$ for some $n\ge1\}$ is a $\g$-submodule, and every
nonzero $\g$-submodule of $G$ contains a nonzero vector of
$P_\nu:=\ker E\cap G\cap Q_\nu$ (apply powers of $E$ and take the last
nonzero image). In the shallow region $-a\le c<b$, $E=e_{-c}$ is injective
on $Q$, so $G=0$; only the deep region carries a nonzero locally nilpotent
part.

\section{The corrected E-action}

\begin{lemma}\label{lem:Eaction}
For $x\in V_0$ define linear maps $A_r:V_0\to V_0$ by
\[
 Ex=\sum_{r\ge0}f_0^r\,A_rx,
\]
a finite sum for each $x$, with $A_0=\pi_0\circ\ad(e_0)$ the old operator
$T$. If $Hx=(\mu+\delta)x$ and $m\ge1$, then in $Q$,
\begin{equation}\label{eq:Eaction}
 E\bigl(f_0^{-m}x\bigr)
 =\sum_{r=0}^{m-1}f_0^{-(m-r)}A_rx
 -m(\mu+\delta+m+1)f_0^{-m-1}x.
\end{equation}
\end{lemma}

\begin{proof}
By the triple identity $Ef_0^{-m}=f_0^{-m}E-mf_0^{-m-1}H-m(m+1)f_0^{-m-2}f_0$
and $Ex=\sum_rf_0^rA_rx$; the term $f_0^{-m}f_0^rA_rx$ survives in $Q$ for
$0\le r<m$ as $f_0^{-(m-r)}A_rx$, and vanishes for $r\ge m$.
\end{proof}

\begin{corollary}\label{cor:counterexample}
The earlier formula retaining only $A_0$ is false. For $X=f_{-1}f_1$ one has
$EXu=(f_1h_{-1}+\lambda_1f_{-1}-2f_0)u$, hence
\[
 E\bigl(f_0^{-2}f_{-1}f_1u\bigr)
 =f_0^{-2}(f_1h_{-1}+\lambda_1f_{-1})u
 -2w_1-2(\mu-1)f_0^{-3}f_{-1}f_1u,
\]
and the $A_0$-only formula misses the nonzero term $-2w_1$ (machine-checked
in \texttt{work/check\_audit\_regressions.py}).
\end{corollary}

\section{The primitive space}

\begin{lemma}\label{lem:pole}
If $0\ne v\in Q_\nu$ and $Ev=0$, then $\nu\in\Z_{\ge0}$, and the highest
pole of $v$ is exactly $\nu+1$.
\end{lemma}

\begin{proof}
Write $v=\sum_{m=1}^{N}f_0^{-m}x_m$ with $x_m\in V_0$. At pole $N+1$ only
the diagonal of \eqref{eq:Eaction} survives, forcing $N=\nu+1$.
\end{proof}

\begin{theorem}\label{thm:P}
A nonzero vector $v=\sum_{m=1}^{\nu+1}f_0^{-m}x_m\in Q_\nu$ is primitive if
and only if its layers satisfy the triangular system
\begin{equation}\label{eq:recursion}
 \sum_{m=k}^{\nu+1}A_{m-k}x_m=(k-1)(\nu-k+2)x_{k-1},
 \qquad 1\le k\le\nu+1,\quad x_0=0,
\end{equation}
and the bottom condition at $k=1$ is the compatibility
$\sum_{m=1}^{\nu+1}A_{m-1}x_m=0$. The primitive space $P_\nu$ is the linear
space of solutions of this system; it is not in general a set of single
monomials, and the old description $P_\nu\cong\{X:T^{\nu+1}(X)u=0\}$ is
false.
\end{theorem}

\begin{proof}
Equate coefficients of $f_0^{-k}$ in $Ev=0$ using \eqref{eq:Eaction}. For
$k=\nu+1$ the right-hand side is $\nu x_\nu=A_0x_{\nu+1}$, so the top
layers are determined downward; the $k=1$ equation is the remaining
compatibility condition. For $\nu=1$ the system reads
$x_1=A_0x_2$ and $(A_0^2+A_1)x_2=0$.
\end{proof}

\begin{corollary}\label{cor:Zcounter}
The old criterion genuinely fails. The element
\[
 Z=h_{-1}^2+4f_{-1}e_{-1}+2h_{-2}
 =h_{-1}^2+2(e_{-1}f_{-1}+f_{-1}e_{-1})
\]
commutes with $E$, $f_0$, and $H$ (it need not be central in $\g$). At
$\mu=-1$, the vector $p=\lambda_1f_0^{-1}u+f_0^{-2}f_1u$ lies in $P_1$, and
$Zp\in P_1$ has layers
$x_1=(\lambda_1Z-4h_{-1})u$, $x_2=(f_1Z-4(\kappa+1)f_{-1})u$ with
$A_0^2x_2=8e_{-1}u\ne0$; the missing term $A_1x_2=-8e_{-1}u$ restores the
compatibility (machine-checked).
\end{corollary}

\begin{corollary}\label{cor:P0}
At $\mu=-2$, $\dim P_0=\infty$: the vectors
$Z^nw_1=f_0^{-1}Z^nu+M$ ($n\ge0$) are primitive and linearly independent,
their leading PBW degrees being $2n$ with leading symbols
$(h_{-1}^2+4f_{-1}e_{-1})^n\ne0$. The earlier claim $P_0=\C w_1$ is
withdrawn.
\end{corollary}

\section{The \texorpdfstring{$\slc$}{sl2}-decomposition and the projections}

\begin{theorem}\label{thm:sl2}
Every vector of $G$ lies in a finite-dimensional
$\slc(E,H,f_0)$-submodule, and
\[
 G\;\cong\;\bigoplus_{\nu\in\Z_{\ge0}}P_\nu\otimes L(\nu)
\]
as transverse $\slc$-modules, where $L(\nu)$ is the $(\nu+1)$-dimensional
irreducible and $P_\nu$ is the multiplicity space. This decomposition is
correct and is retained; it is not a decomposition into $\g$-submodules.
\end{theorem}

\begin{proof}
$E$ and $f_0$ are locally nilpotent on $G$, so each vector generates a
finite-dimensional $\slc$-module; complete reducibility is Weyl's theorem,
and primitive weights are nonnegative integers.
\end{proof}

The principle behind all projections is that the current modes form the
adjoint module $L(2)$ under the transverse $\slc$, and
$L(2)\otimes L(\nu)\cong L(\nu+2)\oplus L(\nu)\oplus L(\nu-2)$ for
$\nu\ge2$.

\begin{proposition}\label{prop:proj}
Let $p\in P_\nu$ have top layer $X$ at pole $\nu+1$.
\begin{enumerate}
\item[(i)] For $i\le A$: $e_ip\in P_{\nu+2}$, and its top layer is
$-2f_i u$ for $\nu=0$ (example: $e_{-1}w_1=f_0^{-1}e_{-1}u-f_0^{-2}h_{-1}u
-2f_0^{-3}f_{-1}u$ at $\mu=-2$), and in general the top pole $\nu+3$ carries
the correction layer of $[e_i,f_0^{-m}]$.
\item[(ii)] For $r<0$ and $\nu\ge1$:
$h_rp+\tfrac{2}{\nu+2}f_0e_rp\in P_\nu$, with top layer
$\tfrac{\nu}{\nu+2}(-h_rX+\tfrac1\nu f_rA_0X)$; for $\nu=0$,
$h_rp=-f_0e_rp$.
\item[(iii)] For $j\ge1$: $E^2f_jp=-2e_jp$, and $f_jp$ has a component in
$P_{\nu+2}$ and one in $P_\nu$.
\item[(iv)] For $\nu\ge2$ and all $j$: the lowering projection
\begin{equation}\label{eq:lowering}
 \ell_j(p)=f_jp-\tfrac1\nu f_0h_jp
 -\tfrac{1}{\nu(\nu+1)}f_0^2e_jp\;\in\;P_{\nu-2},
\end{equation}
the $L(\nu-2)$-component of the tensor product above. It is nonzero: at
$\mu=-2$, for $p_2=e_{-1}w_1\in P_2$,
\[
 \ell_1(p_2)=f_0^{-1}\Bigl(\tfrac13f_1e_{-1}+\tfrac{\lambda_1}{6}h_{-1}
 +\tfrac{2\kappa-2}{3}\Bigr)u+M\in P_0\setminus\{0\}.
\]
\end{enumerate}
\end{proposition}

\begin{proof}
(ii) and (iii) use $[E,h_r]=-2e_r$, $[f_0,e_r]=-h_r$ and $[E,f_j]=h_j$ as
before; (iv) is verified by applying $E$ to the right-hand side of
\eqref{eq:lowering} and using $Ef_0^s=s(\nu-s+1)f_0^{s-1}$ on the block of
$p$. The nonvanishing example is machine-checked.
\end{proof}

\section{Submodule structure: what is withdrawn and what remains open}

\paragraph{The weight tails are not affine submodules.}
The earlier note claimed that
$G^{\ge\nu}:=\bigoplus_{\nu'\ge\nu}P_{\nu'}\otimes L(\nu')$ is a
$\g$-submodule. This is false.
\begin{enumerate}
\item[(i)] Central obstruction: at $\mu=-2$ and $\kappa\ne0$, take
$p=w_1\in L(0)$. Then $h_{\pm1}p=-f_0e_{\pm1}p$ lies in the blocks of top
weight $2$, so both $h_1p$ and $h_{-1}p$ belong to $G^{\ge2}$. Were
$G^{\ge2}$ a submodule, $[h_1,h_{-1}]p=2\kappa p$ would belong to it,
contradicting $p\in L(0)$ and $\kappa\ne0$ (machine-checked).
\item[(ii)] The lowering channel: $\ell_1(e_{-1}w_1)\in P_0\setminus\{0\}$
at $\mu=-2$ shows that modes map $P_2$ into $P_0$, so the tails are not
stable even at $\kappa=0$.
\end{enumerate}
Consequently there are no nonzero trivial affine quotient modules of $Q$ at
level $\kappa\ne0$: $K$ acts by $\kappa$ while all $h$-modes acting by
zero would force $[h_1,h_{-1}]=2K$ to act by zero.

\paragraph{Withdrawn claims.}
The following statements, derived in the earlier version from the invalid
filtration, are withdrawn as proved results: that $G$ is never simple; that
$\operatorname{soc}(G)=0$ and $\operatorname{soc}(Q)=0$; that $G$ has
infinite composition length; that the quotients $G^{\ge\nu}/G^{\ge\nu+1}$
are simple and form the head. Their negations are not proved either: the
internal classification of $G_c$ is open.

\paragraph{What remains proved.}
The boundary isomorphism, the deep-region simplicity criterion
$\mathcal T_{f_c}M_{a,b}$ simple $\iff\lambda_0-c\kappa\notin\Z$ (under
$\lambda_L\ne0$, $\kappa\ne-2$), the unique maximal submodule $G_c$ with
simple head $Q_c/G_c$ and nonsplit extension at integral parameters, the
algebraic central reduction compatibility, the transverse
$\slc$-decomposition of $G_c$, and the corrected recursion of
Theorem~\ref{thm:P} with the projections of Proposition~\ref{prop:proj}.
A correct submodule analysis of $G_c$ must track all three channels
$L(\nu+2)$, $L(\nu)$, $L(\nu-2)$ simultaneously; the old filtration is
abandoned.

\section{Verification status}

The audit \texttt{audits/2026-09-22/} (report, independent checker with 154
exact checks, results, log) is part of this repository. Its four classes of
counterexamples are reproduced by \texttt{work/check\_audit\_regressions.py}
(24 checks): the E-action residual $-2w_1$; the truncated-$T$ residual
$-4e_{-2}u$; the $Z$-family in $P_0$ and the $\mu=-1$ corrected-recursion
example; the obstruction $[h_1,h_{-1}]w_1=2\kappa w_1$; the nonzero
$\ell_1(p_2)$. These are corroborations, not formal proofs, and the open
status of the submodule classification should be reflected in any summary
of this project.

\begin{thebibliography}{9}
\bibitem{FGXZ}
V.~Futorny, X.~Guo, Y.~Xue and K.~Zhao,
\emph{Smooth representations of affine Kac--Moody algebras},
arXiv:2404.03855v2 (2024), Theorem~1.1.
\url{https://arxiv.org/abs/2404.03855v2}.

\bibitem{Audit}
Independent audit of 2026-09-22,
\texttt{audits/2026-09-22/audit\_report\_cn.md} and
\texttt{audits/2026-09-22/independent\_pbw\_audit.py} (154 exact checks).

\bibitem{Unified}
Project working notes: \emph{Boundary and Nonboundary Localization Quotients
of Smooth Affine Modules} (unified note, September 2026) and the addendum
of September 15, 2026.
\end{thebibliography}
\end{document}
